\subsubsection{\label{subsec.sign.pie.moeb-bool}Boolean M\"{o}bius inversion}

This generalization (we will soon see why it is a generalization) is our third
\textquotedblleft Principle of Inclusion and Exclusion\textquotedblright, but
is probably best referred to as the \emph{Boolean M\"{o}bius inversion
formula} (or the \emph{M\"{o}bius inversion formula for the Boolean lattice}):

\begin{theorem}
[Boolean M\"{o}bius inversion]\label{thm.pie.moeb}Let $S$ be a finite set. Let
$A$ be any additive abelian group.

For each subset $I$ of $S$, let $a_{I}$ and $b_{I}$ be two elements of $A$.

Assume that%
\begin{equation}
b_{I}=\sum_{J\subseteq I}a_{J}\ \ \ \ \ \ \ \ \ \ \text{for all }I\subseteq S.
\label{eq.thm.pie.moeb.ass}%
\end{equation}


Then, we also have%
\begin{equation}
a_{I}=\sum_{J\subseteq I}\left(  -1\right)  ^{\left\vert I\setminus
J\right\vert }b_{J}\ \ \ \ \ \ \ \ \ \ \text{for all }I\subseteq S.
\label{eq.thm.pie.moeb.claim}%
\end{equation}

\end{theorem}

\begin{example}
Let $S=\left[  2\right]  =\left\{  1,2\right\}  $. Then, the assumptions of
Theorem \ref{thm.pie.moeb} state that%
\begin{align*}
b_{\varnothing}  &  =a_{\varnothing};\\
b_{\left\{  1\right\}  }  &  =a_{\varnothing}+a_{\left\{  1\right\}  };\\
b_{\left\{  2\right\}  }  &  =a_{\varnothing}+a_{\left\{  2\right\}  };\\
b_{\left\{  1,2\right\}  }  &  =a_{\varnothing}+a_{\left\{  1\right\}
}+a_{\left\{  2\right\}  }+a_{\left\{  1,2\right\}  }.
\end{align*}
The claim of Theorem \ref{thm.pie.moeb} then states that%
\begin{align*}
a_{\varnothing}  &  =b_{\varnothing};\\
a_{\left\{  1\right\}  }  &  =-b_{\varnothing}+b_{\left\{  1\right\}  };\\
a_{\left\{  2\right\}  }  &  =-b_{\varnothing}+b_{\left\{  2\right\}  };\\
a_{\left\{  1,2\right\}  }  &  =b_{\varnothing}-b_{\left\{  1\right\}
}-b_{\left\{  2\right\}  }+b_{\left\{  1,2\right\}  }.
\end{align*}
These four equalities can be verified easily. For instance, let us check the
last of them:%
\begin{align*}
&  \underbrace{b_{\varnothing}}_{=a_{\varnothing}}-\underbrace{b_{\left\{
1\right\}  }}_{=a_{\varnothing}+a_{\left\{  1\right\}  }}%
-\underbrace{b_{\left\{  2\right\}  }}_{=a_{\varnothing}+a_{\left\{
2\right\}  }}+\underbrace{b_{\left\{  1,2\right\}  }}_{=a_{\varnothing
}+a_{\left\{  1\right\}  }+a_{\left\{  2\right\}  }+a_{\left\{  1,2\right\}
}}\\
&  =a_{\varnothing}-\left(  a_{\varnothing}+a_{\left\{  1\right\}  }\right)
-\left(  a_{\varnothing}+a_{\left\{  2\right\}  }\right)  +\left(
a_{\varnothing}+a_{\left\{  1\right\}  }+a_{\left\{  2\right\}  }+a_{\left\{
1,2\right\}  }\right) \\
&  =a_{\left\{  1,2\right\}  }.
\end{align*}

\end{example}

Before we prove Theorem \ref{thm.pie.moeb}, let us show that the weighted PIE
(Theorem \ref{thm.pie.2}) is a particular case of it:

\begin{proof}
[Proof of Theorem \ref{thm.pie.2} using Theorem \ref{thm.pie.moeb}.]Let
$S=\left[  n\right]  $. We note that the map
\begin{align*}
\left\{  \text{subsets of }S\right\}   &  \rightarrow\left\{  \text{subsets of
}S\right\}  ,\\
J  &  \mapsto S\setminus J
\end{align*}
is a bijection. (Indeed, this map is an involution, since each subset $J$ of
$S$ satisfies $S\setminus\left(  S\setminus J\right)  =J$.)

For each $u\in U$, define a subset $\operatorname*{Viol}u$ of $S$ by%
\[
\operatorname*{Viol}u:=\left\{  i\in S\ \mid\ u\notin A_{i}\right\}  .
\]
(In terms of the \textquotedblleft rule-breaking\textquotedblright%
\ interpretation, $\operatorname*{Viol}u$ is the set of all rules that $u$
violates.) Now, for each subset $I$ of $\left[  n\right]  $, we set%
\[
a_{I}:=\sum_{\substack{u\in U;\\\operatorname*{Viol}u=I}}w\left(  u\right)
\ \ \ \ \ \ \ \ \ \ \text{and}\ \ \ \ \ \ \ \ \ \ b_{I}:=\sum_{\substack{u\in
U;\\\operatorname*{Viol}u\subseteq I}}w\left(  u\right)  .
\]
Then, for each subset $I$ of $S$, we have%
\begin{align*}
b_{I}  &  =\sum_{\substack{u\in U;\\\operatorname*{Viol}u\subseteq I}}w\left(
u\right)  =\sum_{J\subseteq I}\ \ \underbrace{\sum_{\substack{u\in
U;\\\operatorname*{Viol}u=J}}w\left(  u\right)  }_{\substack{=a_{J}\\\text{(by
the definition of }a_{J}\text{)}}}\ \ \ \ \ \ \ \ \ \ \left(
\begin{array}
[c]{c}%
\text{here, we have split}\\
\text{the sum according to}\\
\text{the value of }\operatorname*{Viol}u
\end{array}
\right) \\
&  =\sum_{J\subseteq I}a_{J}.
\end{align*}
Thus, we can apply Theorem \ref{thm.pie.moeb}. This gives us%
\begin{equation}
a_{I}=\sum_{J\subseteq I}\left(  -1\right)  ^{\left\vert I\setminus
J\right\vert }b_{J}\ \ \ \ \ \ \ \ \ \ \text{for all }I\subseteq S.
\label{pf.thm.pie.2.4}%
\end{equation}


However, if $J$ is any subset of $S$, then the definition of $b_{J}$ yields%
\begin{equation}
b_{J}=\sum_{\substack{u\in U;\\\operatorname*{Viol}u\subseteq J}}w\left(
u\right)  =\sum_{\substack{u\in U;\\u\in A_{i}\text{ for all }i\in S\setminus
J}}w\left(  u\right)  \label{pf.thm.pie.2.5}%
\end{equation}
(because the elements $u\in U$ that satisfy $\operatorname*{Viol}u\subseteq J$
are precisely the elements $u\in U$ that satisfy $\left(  u\in A_{i}\text{ for
all }i\in S\setminus J\right)  $\ \ \ \ \footnote{\textit{Proof.} Let $u\in
U$. We must prove that the condition \textquotedblleft$\operatorname*{Viol}%
u\subseteq J$\textquotedblright\ is equivalent to \textquotedblleft$u\in
A_{i}$ for all $i\in S\setminus J$\textquotedblright.
\par
We have $\operatorname*{Viol}u=\left\{  i\in S\ \mid\ u\notin A_{i}\right\}  $
(by the definition of $\operatorname*{Viol}u$). Hence, we have the following
chain of equivalences:%
\begin{align*}
\left(  \operatorname*{Viol}u\subseteq J\right)  \  &  \Longleftrightarrow
\ \left(  \left\{  i\in S\ \mid\ u\notin A_{i}\right\}  \subseteq J\right) \\
&  \Longleftrightarrow\ \left(  \text{each }i\in S\text{ satisfying }u\notin
A_{i}\text{ belongs to }J\right) \\
&  \Longleftrightarrow\ \left(  \text{each }i\in S\text{ that does not belong
to }J\text{ must satisfy }u\in A_{i}\right) \\
&  \ \ \ \ \ \ \ \ \ \ \ \ \ \ \ \ \ \ \ \ \left(  \text{by contraposition}%
\right) \\
&  \Longleftrightarrow\ \left(  \text{each }i\in S\setminus J\text{ must
satisfy }u\in A_{i}\right) \\
&  \Longleftrightarrow\ \left(  u\in A_{i}\text{ for all }i\in S\setminus
J\right)  .
\end{align*}
Hence, the condition \textquotedblleft$\operatorname*{Viol}u\subseteq
J$\textquotedblright\ is equivalent to \textquotedblleft$u\in A_{i}$ for all
$i\in S\setminus J$\textquotedblright.}).

Now, applying (\ref{pf.thm.pie.2.4}) to $I=S$, we obtain%
\begin{align*}
a_{S}  &  =\sum_{J\subseteq S}\left(  -1\right)  ^{\left\vert S\setminus
J\right\vert }b_{J}=\sum_{J\subseteq S}\left(  -1\right)  ^{\left\vert
S\setminus J\right\vert }\sum_{\substack{u\in U;\\u\in A_{i}\text{ for all
}i\in S\setminus J}}w\left(  u\right)  \ \ \ \ \ \ \ \ \ \ \left(  \text{by
(\ref{pf.thm.pie.2.5})}\right) \\
&  =\sum_{I\subseteq S}\left(  -1\right)  ^{\left\vert I\right\vert }%
\sum_{\substack{u\in U;\\u\in A_{i}\text{ for all }i\in I}}w\left(  u\right)
\\
&  \ \ \ \ \ \ \ \ \ \ \ \ \ \ \ \ \ \ \ \ \left(
\begin{array}
[c]{c}%
\text{here, we have substituted }I\text{ for }S\setminus J\text{ in the sum,
since}\\
\text{the map }\left\{  \text{subsets of }S\right\}  \rightarrow\left\{
\text{subsets of }S\right\}  ,\ J\mapsto S\setminus J\\
\text{is a bijection}%
\end{array}
\right) \\
&  =\sum_{I\subseteq\left[  n\right]  }\left(  -1\right)  ^{\left\vert
I\right\vert }\sum_{\substack{u\in U;\\u\in A_{i}\text{ for all }i\in
I}}w\left(  u\right)  \ \ \ \ \ \ \ \ \ \ \left(  \text{since }S=\left[
n\right]  \right)  .
\end{align*}


On the other hand, the definition of $a_{S}$ yields%
\[
a_{S}=\sum_{\substack{u\in U;\\\operatorname*{Viol}u=S}}w\left(  u\right)
=\sum_{\substack{u\in U;\\u\notin A_{i}\text{ for all }i\in\left[  n\right]
}}w\left(  u\right)
\]
(since the elements $u\in U$ that satisfy $\operatorname*{Viol}u=S$ are
precisely the elements $u\in U$ that satisfy $\left(  u\notin A_{i}\text{ for
all }i\in\left[  n\right]  \right)  $\ \ \ \ \footnote{\textit{Proof.} Let
$u\in U$. We must prove that the condition \textquotedblleft%
$\operatorname*{Viol}u=S$\textquotedblright\ is equivalent to
\textquotedblleft$u\notin A_{i}$ for all $i\in\left[  n\right]  $%
\textquotedblright.
\par
We have $\operatorname*{Viol}u=\left\{  i\in S\ \mid\ u\notin A_{i}\right\}  $
(by the definition of $\operatorname*{Viol}u$). Hence, we have the following
chain of equivalences:%
\begin{align*}
\left(  \operatorname*{Viol}u=S\right)  \  &  \Longleftrightarrow\ \left(
\left\{  i\in S\ \mid\ u\notin A_{i}\right\}  =S\right) \\
&  \Longleftrightarrow\ \left(  \text{each }i\in S\text{ satisfies }u\notin
A_{i}\right) \\
&  \Longleftrightarrow\ \left(  u\notin A_{i}\text{ for all }i\in S\right) \\
&  \Longleftrightarrow\ \left(  u\notin A_{i}\text{ for all }i\in\left[
n\right]  \right)  \ \ \ \ \ \ \ \ \ \ \left(  \text{since }S=\left[
n\right]  \right)  .
\end{align*}
Hence, the condition \textquotedblleft$\operatorname*{Viol}u=S$%
\textquotedblright\ is equivalent to \textquotedblleft$u\notin A_{i}$ for all
$i\in\left[  n\right]  $\textquotedblright.}).

Comparing these two equalities, we obtain%
\[
\sum_{\substack{u\in U;\\u\notin A_{i}\text{ for all }i\in\left[  n\right]
}}w\left(  u\right)  =\sum_{I\subseteq\left[  n\right]  }\left(  -1\right)
^{\left\vert I\right\vert }\sum_{\substack{u\in U;\\u\in A_{i}\text{ for all
}i\in I}}w\left(  u\right)  .
\]
Thus, Theorem \ref{thm.pie.2} has been proved, assuming Theorem
\ref{thm.pie.moeb}.
\end{proof}

It remains to prove the latter:

\begin{proof}
[Proof of Theorem \ref{thm.pie.moeb}.]Fix a subset $Q$ of $S$. We shall prove
that%
\[
a_{Q}=\sum_{I\subseteq Q}\left(  -1\right)  ^{\left\vert Q\setminus
I\right\vert }b_{I}.
\]


We begin by rewriting the right hand side:%
\begin{align}
\sum_{I\subseteq Q}\left(  -1\right)  ^{\left\vert Q\setminus I\right\vert
}\underbrace{b_{I}}_{\substack{=\sum_{J\subseteq I}a_{J}\\\text{(by
(\ref{eq.thm.pie.moeb.ass}))}}}  &  =\sum_{I\subseteq Q}\left(  -1\right)
^{\left\vert Q\setminus I\right\vert }\underbrace{\sum_{J\subseteq I}a_{J}%
}_{=\sum_{P\subseteq I}a_{P}}=\sum_{I\subseteq Q}\left(  -1\right)
^{\left\vert Q\setminus I\right\vert }\sum_{P\subseteq I}a_{P}\nonumber\\
&  =\sum_{I\subseteq Q}\ \ \sum_{P\subseteq I}\left(  -1\right)  ^{\left\vert
Q\setminus I\right\vert }a_{P}. \label{pf.thm.pie.moeb.1}%
\end{align}


The two summation signs \textquotedblleft$\sum_{I\subseteq Q}\ \ \sum
_{P\subseteq I}$\textquotedblright\ on the right hand side of this equality
result in a sum over all pairs $\left(  I,P\right)  $ of subsets of $Q$
satisfying $P\subseteq I\subseteq Q$. The same result can be obtained by the
two summation signs \textquotedblleft$\sum_{P\subseteq Q}\ \ \sum
_{\substack{I\subseteq Q;\\P\subseteq I}}$\textquotedblright\ (indeed, the
only difference between \textquotedblleft$\sum_{I\subseteq Q}\ \ \sum
_{P\subseteq I}$\textquotedblright\ and \textquotedblleft$\sum_{P\subseteq
Q}\ \ \sum_{\substack{I\subseteq Q;\\P\subseteq I}}$\textquotedblright\ is the
order in which the two subsets $I$ and $P$ are chosen). Thus, we can replace
\textquotedblleft$\sum_{I\subseteq Q}\ \ \sum_{P\subseteq I}$%
\textquotedblright\ by \textquotedblleft$\sum_{P\subseteq Q}\ \ \sum
_{\substack{I\subseteq Q;\\P\subseteq I}}$\textquotedblright\ on the right
hand side of (\ref{pf.thm.pie.moeb.1}). Hence, (\ref{pf.thm.pie.moeb.1})
rewrites as follows:%
\begin{align}
\sum_{I\subseteq Q}\left(  -1\right)  ^{\left\vert Q\setminus I\right\vert
}b_{I}  &  =\sum_{P\subseteq Q}\ \ \sum_{\substack{I\subseteq Q;\\P\subseteq
I}}\left(  -1\right)  ^{\left\vert Q\setminus I\right\vert }a_{P}\nonumber\\
&  =\sum_{P\subseteq Q}\left(  \sum_{\substack{I\subseteq Q;\\P\subseteq
I}}\left(  -1\right)  ^{\left\vert Q\setminus I\right\vert }\right)  a_{P}.
\label{pf.thm.pie.moeb.2}%
\end{align}


We want to prove that this equals $a_{Q}$. Since the $a_{P}$'s are arbitrary
elements of an abelian group, the only way this can possibly be achieved is by
showing that the sum on the right hand side simplifies to $a_{Q}$
\textbf{formally} -- i.e., that the coefficient $\sum_{\substack{I\subseteq
Q;\\P\subseteq I}}\left(  -1\right)  ^{\left\vert Q\setminus I\right\vert }$
in front of $a_{P}$ is $0$ whenever $P\neq Q$, and is $1$ whenever $P=Q$.
Thus, we now set out to prove this. Using Definition \ref{def.iverson}, we can
restate this goal as follows: We want to prove that every subset $P$ of $Q$
satisfies%
\begin{equation}
\sum_{\substack{I\subseteq Q;\\P\subseteq I}}\left(  -1\right)  ^{\left\vert
Q\setminus I\right\vert }=\left[  P=Q\right]  . \label{pf.thm.pie.moeb.4}%
\end{equation}


We shall prove this soon (in Lemma \ref{lem.pie.two-sets-altsum} \textbf{(b)}
below). For now, let us explain how the proof of Theorem \ref{thm.pie.moeb}
can be completed if (\ref{pf.thm.pie.moeb.4}) is known to be true. Indeed,
(\ref{pf.thm.pie.moeb.2}) becomes%
\begin{align*}
\sum_{I\subseteq Q}\left(  -1\right)  ^{\left\vert Q\setminus I\right\vert
}b_{I}  &  =\sum_{P\subseteq Q}\underbrace{\left(  \sum_{\substack{I\subseteq
Q;\\P\subseteq I}}\left(  -1\right)  ^{\left\vert Q\setminus I\right\vert
}\right)  }_{\substack{=\left[  P=Q\right]  \\\text{(by
(\ref{pf.thm.pie.moeb.4}))}}}a_{P}=\sum_{P\subseteq Q}\left[  P=Q\right]
a_{P}\\
&  =\underbrace{\left[  Q=Q\right]  }_{\substack{=1\\\text{(since }%
Q=Q\text{)}}}a_{Q}+\sum_{\substack{P\subseteq Q;\\P\neq Q}}\underbrace{\left[
P=Q\right]  }_{\substack{=0\\\text{(since }P\neq Q\text{)}}}a_{P}\\
&  \ \ \ \ \ \ \ \ \ \ \ \ \ \ \ \ \ \ \ \ \left(
\begin{array}
[c]{c}%
\text{here, we have split off the addend for }P=Q\\
\text{from the sum (since }Q\subseteq Q\text{)}%
\end{array}
\right) \\
&  =a_{Q}+\underbrace{\sum_{\substack{P\subseteq Q;\\P\neq Q}}0a_{P}}%
_{=0}=a_{Q}.
\end{align*}
In other words, $a_{Q}=\sum_{I\subseteq Q}\left(  -1\right)  ^{\left\vert
Q\setminus I\right\vert }b_{I}$.

Forget that we fixed $Q$. We thus have shown that%
\[
a_{Q}=\sum_{I\subseteq Q}\left(  -1\right)  ^{\left\vert Q\setminus
I\right\vert }b_{I}\ \ \ \ \ \ \ \ \ \ \text{for all }Q\subseteq S.
\]
Renaming the indices $Q$ and $I$ as $I$ and $J$ in this statement, we obtain
the following:%
\[
a_{I}=\sum_{J\subseteq I}\left(  -1\right)  ^{\left\vert I\setminus
J\right\vert }b_{J}\ \ \ \ \ \ \ \ \ \ \text{for all }I\subseteq S.
\]
Thus, Theorem \ref{thm.pie.moeb} is proved (assuming that
(\ref{pf.thm.pie.moeb.4}) is known to be true).
\end{proof}

It now remains to prove (\ref{pf.thm.pie.moeb.4}). We shall do this as part of
the following lemma:

\begin{lemma}
\label{lem.pie.two-sets-altsum}Let $Q$ be a finite set. Let $P$ be a subset of
$Q$. Then: \medskip

\textbf{(a)} We have%
\begin{equation}
\sum_{\substack{I\subseteq Q;\\P\subseteq I}}\left(  -1\right)  ^{\left\vert
I\right\vert }=\left(  -1\right)  ^{\left\vert P\right\vert }\left[
P=Q\right]  . \label{eq.lem.pie.two-sets-altsum.a}%
\end{equation}


\textbf{(b)} We have%
\begin{equation}
\sum_{\substack{I\subseteq Q;\\P\subseteq I}}\left(  -1\right)  ^{\left\vert
Q\setminus I\right\vert }=\left[  P=Q\right]  .
\label{eq.lem.pie.two-sets-altsum.b}%
\end{equation}

\end{lemma}

As promised, Lemma \ref{lem.pie.two-sets-altsum} \textbf{(b)} (once proved)
will yield (\ref{pf.thm.pie.moeb.4}) and thus will complete our above proof of
Theorem \ref{thm.pie.moeb} (and, with it, the proofs of Theorem
\ref{thm.pie.2} and Theorem \ref{thm.pie.1}).

\begin{proof}
[Proof of Lemma \ref{lem.pie.two-sets-altsum}.]\textbf{(a)} There are many
ways to prove this (in particular, a simple one using the binomial theorem --
do you see it?); but staying true to the spirit of this chapter, we pick one
using a sign-reversing involution. (A variant of this proof can be found in
\cite[solution to Exercise 2.9.1]{19fco}.\footnote{Our sets $Q$ and $P$ are
called $S$ and $T$ in \cite[solution to Exercise 2.9.1]{19fco}.})

We must prove the equality (\ref{eq.lem.pie.two-sets-altsum.a}). If $P=Q$,
then this equality is easily seen to hold\footnote{\textit{Proof.} Assume that
$P=Q$. Then, the only subset $I$ of $Q$ that satisfies $P\subseteq I$ is the
set $Q$ itself (since any such subset $I$ has to satisfy both $Q=P\subseteq I$
and $I\subseteq Q$, which in combination entail $I=Q$). Thus, the sum
$\sum_{\substack{I\subseteq Q;\\P\subseteq I}}\left(  -1\right)  ^{\left\vert
I\right\vert }$ has only one addend, namely the addend for $I=Q$.
Consequently, this sum simplifies as follows:%
\begin{equation}
\sum_{\substack{I\subseteq Q;\\P\subseteq I}}\left(  -1\right)  ^{\left\vert
I\right\vert }=\left(  -1\right)  ^{\left\vert Q\right\vert }.
\label{pf.lem.pie.two-sets-altsum.a.fn1.1}%
\end{equation}
On the other hand, from $P=Q$, we obtain%
\[
\left(  -1\right)  ^{\left\vert P\right\vert }\left[  P=Q\right]  =\left(
-1\right)  ^{\left\vert Q\right\vert }\underbrace{\left[  Q=Q\right]
}_{\substack{=1\\\text{(since }Q=Q\text{)}}}=\left(  -1\right)  ^{\left\vert
Q\right\vert }.
\]
Comparing this with (\ref{pf.lem.pie.two-sets-altsum.a.fn1.1}), we obtain
$\sum_{\substack{I\subseteq Q;\\P\subseteq I}}\left(  -1\right)  ^{\left\vert
I\right\vert }=\left(  -1\right)  ^{\left\vert P\right\vert }\left[
P=Q\right]  $. Thus, we have shown that (\ref{eq.lem.pie.two-sets-altsum.a})
holds under the assumption that $P=Q$.}. Hence, for the rest of this proof, we
WLOG assume that $P\neq Q$. Thus, $\left[  P=Q\right]  =0$.

Now, $P$ is a \textbf{proper} subset of $Q$ (since $P$ is a subset of $Q$ and
satisfies $P\neq Q$). Hence, there exists some $q\in Q$ such that $q\notin P$.
Fix such a $q$.

Let
\[
\mathcal{A}:=\left\{  I\subseteq Q\ \mid\ P\subseteq I\right\}  ,
\]
and let
\[
\operatorname*{sign}I:=\left(  -1\right)  ^{\left\vert I\right\vert
}\ \ \ \ \ \ \ \ \ \ \text{for each }I\in\mathcal{A}.
\]
Then,%
\begin{equation}
\sum_{I\in\mathcal{A}}\operatorname*{sign}I=\sum_{\substack{I\subseteq
Q;\\P\subseteq I}}\left(  -1\right)  ^{\left\vert I\right\vert }.
\label{pf.lem.pie.two-sets-altsum.a.3}%
\end{equation}
We shall now construct an involution $f:\mathcal{A}\rightarrow\mathcal{A}$ on
the set $\mathcal{A}$; this will allow us to apply Lemma
\ref{lem.sign.cancel2} (to $\mathcal{X}=\mathcal{A}$), and easily conclude
that $\sum_{I\in\mathcal{A}}\operatorname*{sign}I=0$ (see below for the details).

Indeed, if $I$ is a subset of $Q$, then $I\cup\left\{  q\right\}  $ is also a
subset of $Q$ (because $q\in Q$). Thus, if $I\in\mathcal{A}$, then
$I\cup\left\{  q\right\}  \in\mathcal{A}$ (because $P\subseteq I$ implies
$P\subseteq I\subseteq I\cup\left\{  q\right\}  $).

On the other hand, if $I$ is a set satisfying $P\subseteq I$, then the set
$I\setminus\left\{  q\right\}  $ also satisfies $P\subseteq I\setminus\left\{
q\right\}  $ (since $q\notin P$). Thus, if $I\in\mathcal{A}$, then
$I\setminus\left\{  q\right\}  \in\mathcal{A}$ (since $I\subseteq Q$ implies
$I\setminus\left\{  q\right\}  \subseteq I\subseteq Q$).

Now, we define a map $f:\mathcal{A}\rightarrow\mathcal{A}$ by
setting\footnote{Here, the notation $X\bigtriangleup Y$ means the symmetric
difference $\left(  X\cup Y\right)  \setminus\left(  X\cap Y\right)  $ of two
sets $X$ and $Y$ (as in Subsection \ref{subsec.gf.defs.commrings}).}%
\[
f\left(  I\right)  :=I\bigtriangleup\left\{  q\right\}  =%
\begin{cases}
I\setminus\left\{  q\right\}  , & \text{if }q\in I;\\
I\cup\left\{  q\right\}  , & \text{if }q\notin I
\end{cases}
\ \ \ \ \ \ \ \ \ \ \text{for each }I\in\mathcal{A}.
\]
This map $f$ is well-defined, because (as we have just shown in the two
paragraphs above) every $I\in\mathcal{A}$ satisfies $I\cup\left\{  q\right\}
\in\mathcal{A}$ and $I\setminus\left\{  q\right\}  \in\mathcal{A}$. Moreover,
this map $f$ is an involution\footnote{Indeed, the map $f$ merely removes $q$
from a set $I$ if $q$ is contained in $I$, and inserts it into $I$ otherwise;
but this is clearly an operation that undoes itself when performed a second
time.}. This involution $f$ has no fixed points (because if $I\in\mathcal{A}$,
then $f\left(  I\right)  =I\bigtriangleup\left\{  q\right\}  \neq I$).
Furthermore, if $I\in\mathcal{A}$, then the set $f\left(  I\right)
=I\bigtriangleup\left\{  q\right\}  $ differs from $I$ in exactly one element
(namely, $q$), and thus satisfies $\left\vert f\left(  I\right)  \right\vert
=\left\vert I\right\vert \pm1$, so that%
\[
\left(  -1\right)  ^{\left\vert f\left(  I\right)  \right\vert }=-\left(
-1\right)  ^{\left\vert I\right\vert },
\]
or, equivalently,%
\[
\operatorname*{sign}\left(  f\left(  I\right)  \right)  =-\operatorname*{sign}%
I
\]
(since the definition of $\operatorname*{sign}I$ yields $\operatorname*{sign}%
I=\left(  -1\right)  ^{\left\vert I\right\vert }$, and similarly
$\operatorname*{sign}\left(  f\left(  I\right)  \right)  =\left(  -1\right)
^{\left\vert f\left(  I\right)  \right\vert }$). Thus, Lemma
\ref{lem.sign.cancel2} (applied to $\mathcal{X}=\mathcal{A}$) shows that%
\begin{align*}
\sum_{I\in\mathcal{A}}\operatorname*{sign}I  &  =\sum_{I\in\mathcal{A}%
\setminus\mathcal{A}}\operatorname*{sign}I=\left(  \text{empty sum}\right)
\ \ \ \ \ \ \ \ \ \ \left(  \text{since }\mathcal{A}\setminus\mathcal{A}%
=\varnothing\right) \\
&  =0.
\end{align*}
Comparing this with (\ref{pf.lem.pie.two-sets-altsum.a.3}), we find%
\[
\sum_{\substack{I\subseteq Q;\\P\subseteq I}}\left(  -1\right)  ^{\left\vert
I\right\vert }=0=\left(  -1\right)  ^{\left\vert P\right\vert }\left[
P=Q\right]  \ \ \ \ \ \ \ \ \ \ \left(  \text{since }\left(  -1\right)
^{\left\vert P\right\vert }\underbrace{\left[  P=Q\right]  }_{=0}=0\right)  .
\]
Thus, (\ref{eq.lem.pie.two-sets-altsum.a}) is proved. This proves Lemma
\ref{lem.pie.two-sets-altsum} \textbf{(a)}. \medskip

\textbf{(b)} If $I$ is any subset of $Q$, then $\left\vert Q\setminus
I\right\vert =\left\vert Q\right\vert -\left\vert I\right\vert \equiv
\left\vert Q\right\vert +\left\vert I\right\vert \operatorname{mod}2$ and thus%
\[
\left(  -1\right)  ^{\left\vert Q\setminus I\right\vert }=\left(  -1\right)
^{\left\vert Q\right\vert +\left\vert I\right\vert }=\left(  -1\right)
^{\left\vert Q\right\vert }\left(  -1\right)  ^{\left\vert I\right\vert }.
\]
Hence,%
\begin{align*}
\sum_{\substack{I\subseteq Q;\\P\subseteq I}}\underbrace{\left(  -1\right)
^{\left\vert Q\setminus I\right\vert }}_{=\left(  -1\right)  ^{\left\vert
Q\right\vert }\left(  -1\right)  ^{\left\vert I\right\vert }}  &
=\sum_{\substack{I\subseteq Q;\\P\subseteq I}}\left(  -1\right)  ^{\left\vert
Q\right\vert }\left(  -1\right)  ^{\left\vert I\right\vert }=\left(
-1\right)  ^{\left\vert Q\right\vert }\underbrace{\sum_{\substack{I\subseteq
Q;\\P\subseteq I}}\left(  -1\right)  ^{\left\vert I\right\vert }%
}_{\substack{=\left(  -1\right)  ^{\left\vert P\right\vert }\left[
P=Q\right]  \\\text{(by Lemma \ref{lem.pie.two-sets-altsum} \textbf{(a)})}}}\\
&  =\left(  -1\right)  ^{\left\vert Q\right\vert }\left(  -1\right)
^{\left\vert P\right\vert }\left[  P=Q\right]  .
\end{align*}
However, it is easy to see that%
\begin{equation}
\left(  -1\right)  ^{\left\vert Q\right\vert }\left(  -1\right)  ^{\left\vert
P\right\vert }\left[  P=Q\right]  =\left[  P=Q\right]
\label{pf.lem.pie.two-sets-altsum.b.1}%
\end{equation}
\footnote{\textit{Proof of (\ref{pf.lem.pie.two-sets-altsum.b.1}):} If $P\neq
Q$, then the equality (\ref{pf.lem.pie.two-sets-altsum.b.1}) boils down to
$\left(  -1\right)  ^{\left\vert Q\right\vert }\left(  -1\right)  ^{\left\vert
P\right\vert }\cdot0=0$ (since $P\neq Q$ entails $\left[  P=Q\right]  =0$),
which is obviously true. Hence, (\ref{pf.lem.pie.two-sets-altsum.b.1}) is
proved if $P\neq Q$. Thus, for the rest of this proof, we WLOG assume that
$P=Q$. Hence, $\left(  -1\right)  ^{\left\vert Q\right\vert }\left(
-1\right)  ^{\left\vert P\right\vert }=\left(  -1\right)  ^{\left\vert
Q\right\vert }\left(  -1\right)  ^{\left\vert Q\right\vert }=\left(
-1\right)  ^{\left\vert Q\right\vert +\left\vert Q\right\vert }=1$ (since
$\left\vert Q\right\vert +\left\vert Q\right\vert =2\left\vert Q\right\vert $
is even). Therefore, $\underbrace{\left(  -1\right)  ^{\left\vert Q\right\vert
}\left(  -1\right)  ^{\left\vert P\right\vert }}_{=1}\left[  P=Q\right]
=\left[  P=Q\right]  $. This proves (\ref{pf.lem.pie.two-sets-altsum.b.1}).}.
Thus,%
\[
\sum_{\substack{I\subseteq Q;\\P\subseteq I}}\left(  -1\right)  ^{\left\vert
Q\setminus I\right\vert }=\left(  -1\right)  ^{\left\vert Q\right\vert
}\left(  -1\right)  ^{\left\vert P\right\vert }\left[  P=Q\right]  =\left[
P=Q\right]  .
\]
This proves Lemma \ref{lem.pie.two-sets-altsum} \textbf{(b)}.
\end{proof}

As said above, this completes the proofs of Theorem \ref{thm.pie.moeb}, of
Theorem \ref{thm.pie.2} and of Theorem \ref{thm.pie.1}.

While Theorem \ref{thm.pie.moeb} has played the part of the ultimate
generalization to us, it can be generalized further. Indeed, it is merely a
particular case of \emph{M\"{o}bius inversion for arbitrary posets} (see,
e.g., \cite[Proposition 3.7.1]{Stanley-EC1} or \cite[Theorem 2.3.1]{Martin21}
or \cite[Theorem 5.5.5]{Sagan19} or \cite[Theorem 6.10]{Sam21} or
\cite[Theorem 14.6.4]{Wagner20}).
